% 【本文件写什么】子模块 mm-impl，\subsection 级聚合
% - 职责摘要：页表实现：Sv39、LoongArch、dummy 与 common 共享代码。
% - 子域聚合 lib.rs：os/components/wateros-mm/mm-impl/src/lib.rs
% - 如何重导出 api-v0、feature 下挂载哪些 impl
% - 被谁依赖，syscall、vfs、main 启动链等
% - 短综述 + 下方 \input 展开 api 与 impl 叶子
% 【事实来源】
% - os/components/wateros-mm/mm-impl/
% - docs/exports/ 中与 wateros-mm 相关的 public-api / features 章节

\subsubsection{impl-common}

\code{wateros-mm-impl-common}位于\code{mm-api}之下，不对外暴露稳定契约。各 arch \code{mm-impl}共享 ELF 装载辅助、mmap/mremap 公共路径、按需零页 loader 等逻辑。

\code{ZeroAnonLoader}实现\code{DemandPageLoader}：匿名惰性 mmap 缺页时保留预先清零页，避免 glibc pthread 栈等大区间饥渴分配耗尽帧池。\code{fill\_elf\_load\_page}按\code{PT\_LOAD}段边界从文件偏移填充单页。\code{ElfSegmentLoadParams}封装 execve lazy map 的段参数。\code{executable}侧三次读稳定、shebang 解析等亦集中于此 crate。

语义边界仍以\code{mm-api} trait 为准；更换 loader 策略或 bring-up 假设时只改\code{impl-common}与各 arch \code{kernel\_elf}，不修改 syscall 契约。

\begin{rustcode}
pub struct ZeroAnonLoader;
impl DemandPageLoader for ZeroAnonLoader {
    fn duplicate_box(&self) -> MmResult<Box<dyn DemandPageLoader>> {
        Ok(Box::new(ZeroAnonLoader))
    }
    fn load_page(&mut self, _file_offset: usize, _dst: &mut [u8]) -> MmResult<()> {
        Ok(())
    }
}
\end{rustcode}
